\documentclass[reqno]{amsart}
\usepackage[T1]{fontenc}
\usepackage{amsmath,amssymb,amsthm}
\usepackage{xr-hyper}
\usepackage[hidelinks]{hyperref}
\externaldocument[sm-]{20261125T000000--slope-method__workbook}
\externaldocument[gf-]{20261125T005000--galois-fourier-and-averaging__workbook}
\externaldocument[ip-]{20261124T214600--the-isotropic-plane__workbook}
\externaldocument[hn-]{20261125T013500--hecke-nondegeneracy-on-the-tree__workbook}

% Chapter number; \ref{chap} gives it.
\newcounter{chap}
% Chapter numbers, generated from 20261125T012000--grh-workbook__guide.org
% by scripts/make-chapters.py.  Do not edit: change the order of the
% index and rerun the script (./preview-site.sh and CI do this).
%
% Usage in a sheet, after \newcounter{chap}:
%   % Chapter numbers, generated from 20261125T012000--grh-workbook__guide.org
% by scripts/make-chapters.py.  Do not edit: change the order of the
% index and rerun the script (./preview-site.sh and CI do this).
%
% Usage in a sheet, after \newcounter{chap}:
%   % Chapter numbers, generated from 20261125T012000--grh-workbook__guide.org
% by scripts/make-chapters.py.  Do not edit: change the order of the
% index and rerun the script (./preview-site.sh and CI do this).
%
% Usage in a sheet, after \newcounter{chap}:
%   \input{chapters}\chapnumber{<sheet ID>}
\makeatletter
\providecommand\chapters@set[2]{\@namedef{chapters@#1}{#2}}
\chapters@set{20261125T001000}{1}% zeros-to-ramp-sums
\chapters@set{20261125T001100}{2}% inert-primes-and-estimate-b
\chapters@set{20261125T005000}{3}% galois-fourier-and-averaging
\chapters@set{20261125T005500}{4}% columns-from-a-sphere
\chapters@set{20261125T000000}{5}% slope-method
\chapters@set{20261126T030100}{6}% a-twist-for-each-prime
\chapters@set{20261125T003000}{7}% the-chain-and-the-square-root
\chapters@set{20261125T010000}{8}% rank-lemma-and-zero-estimate
\chapters@set{20261125T010500}{9}% poincare-series-and-near-orthogonality
\chapters@set{20261124T223000}{10}% elliptic-curves-for-the-proof
\chapters@set{20261125T009500}{11}% fock-model-and-theta-intertwiners
\chapters@set{20261124T223100}{12}% complex-multiplication-and-inert-reduction
\chapters@set{20261125T008000}{13}% one-cm-determinant
\chapters@set{20261124T214500}{14}% supersingular-theta-divisibility
\chapters@set{20261124T214600}{15}% the-isotropic-plane
\chapters@set{20261125T013700}{16}% pairings-on-a-formal-group-of-height-two
\chapters@set{20261125T013500}{17}% hecke-nondegeneracy-on-the-tree
\chapters@set{20261126T030200}{18}% the-bound-is-attained
\chapters@set{20261125T009000}{19}% kronecker-limit-and-chowla-selberg
\chapters@set{20261125T002000}{20}% definite-theta-lifts
\chapters@set{20261125T015500}{21}% siegel-weil-in-the-smallest-case
\chapters@set{20261125T002100}{22}% definite-theta-lifts-and-norms
\chapters@set{20261125T016500}{23}% hecke-operators-and-theta-lifts
\chapters@set{20261125T016700}{24}% divisibility-along-the-chain
\chapters@set{20261125T004000}{25}% denominators-and-sha
\chapters@set{20261125T004100}{26}% tate-twists-and-sha
\chapters@set{20261126T020430}{27}% the-minus-part-by-gauss-sums
\chapters@set{20261125T014100}{28}% congruence-modules
\chapters@set{20261125T014500}{29}% eisenstein-congruences-and-selmer-classes
\chapters@set{20261125T014700}{30}% pseudocharacters-and-the-length-bound
\chapters@set{20261125T011000}{31}% galois-representations-black-box
\chapters@set{20261126T030300}{32}% the-data
\chapters@set{20261126T020210}{33}% bessel-poincare-series-and-their-gram-operators
\chapters@set{20261126T030400}{34}% the-seesaw-identity
\chapters@set{20261126T030500}{35}% divisibility-of-the-composite-lift
\chapters@set{20261126T030600}{36}% theorem-a-short-forms
\chapters@set{20261126T030700}{37}% the-determinant-of-a-chain
\chapters@set{20261126T030800}{38}% the-cost-of-one-lift
\chapters@set{20261126T030900}{39}% denominators-of-projections-to-lifts
\chapters@set{20261126T031000}{40}% twists-with-few-congruences
\chapters@set{20261126T031100}{41}% theorem-c-and-the-main-inequality
\providecommand\chapnumber[1]{%
  \@ifundefined{chapters@#1}%
    {\PackageError{chapters}{Sheet #1 is not listed in the index}%
      {List it in 20261125T012000--grh-workbook__guide.org\MessageBreak
       and run scripts/make-chapters.py.}}%
    {\edef\thechap{\@nameuse{chapters@#1}}}}
\makeatother
\chapnumber{<sheet ID>}
\makeatletter
\providecommand\chapters@set[2]{\@namedef{chapters@#1}{#2}}
\chapters@set{20261125T001000}{1}% zeros-to-ramp-sums
\chapters@set{20261125T001100}{2}% inert-primes-and-estimate-b
\chapters@set{20261125T005000}{3}% galois-fourier-and-averaging
\chapters@set{20261125T005500}{4}% columns-from-a-sphere
\chapters@set{20261125T000000}{5}% slope-method
\chapters@set{20261126T030100}{6}% a-twist-for-each-prime
\chapters@set{20261125T003000}{7}% the-chain-and-the-square-root
\chapters@set{20261125T010000}{8}% rank-lemma-and-zero-estimate
\chapters@set{20261125T010500}{9}% poincare-series-and-near-orthogonality
\chapters@set{20261124T223000}{10}% elliptic-curves-for-the-proof
\chapters@set{20261125T009500}{11}% fock-model-and-theta-intertwiners
\chapters@set{20261124T223100}{12}% complex-multiplication-and-inert-reduction
\chapters@set{20261125T008000}{13}% one-cm-determinant
\chapters@set{20261124T214500}{14}% supersingular-theta-divisibility
\chapters@set{20261124T214600}{15}% the-isotropic-plane
\chapters@set{20261125T013700}{16}% pairings-on-a-formal-group-of-height-two
\chapters@set{20261125T013500}{17}% hecke-nondegeneracy-on-the-tree
\chapters@set{20261126T030200}{18}% the-bound-is-attained
\chapters@set{20261125T009000}{19}% kronecker-limit-and-chowla-selberg
\chapters@set{20261125T002000}{20}% definite-theta-lifts
\chapters@set{20261125T015500}{21}% siegel-weil-in-the-smallest-case
\chapters@set{20261125T002100}{22}% definite-theta-lifts-and-norms
\chapters@set{20261125T016500}{23}% hecke-operators-and-theta-lifts
\chapters@set{20261125T016700}{24}% divisibility-along-the-chain
\chapters@set{20261125T004000}{25}% denominators-and-sha
\chapters@set{20261125T004100}{26}% tate-twists-and-sha
\chapters@set{20261126T020430}{27}% the-minus-part-by-gauss-sums
\chapters@set{20261125T014100}{28}% congruence-modules
\chapters@set{20261125T014500}{29}% eisenstein-congruences-and-selmer-classes
\chapters@set{20261125T014700}{30}% pseudocharacters-and-the-length-bound
\chapters@set{20261125T011000}{31}% galois-representations-black-box
\chapters@set{20261126T030300}{32}% the-data
\chapters@set{20261126T020210}{33}% bessel-poincare-series-and-their-gram-operators
\chapters@set{20261126T030400}{34}% the-seesaw-identity
\chapters@set{20261126T030500}{35}% divisibility-of-the-composite-lift
\chapters@set{20261126T030600}{36}% theorem-a-short-forms
\chapters@set{20261126T030700}{37}% the-determinant-of-a-chain
\chapters@set{20261126T030800}{38}% the-cost-of-one-lift
\chapters@set{20261126T030900}{39}% denominators-of-projections-to-lifts
\chapters@set{20261126T031000}{40}% twists-with-few-congruences
\chapters@set{20261126T031100}{41}% theorem-c-and-the-main-inequality
\providecommand\chapnumber[1]{%
  \@ifundefined{chapters@#1}%
    {\PackageError{chapters}{Sheet #1 is not listed in the index}%
      {List it in 20261125T012000--grh-workbook__guide.org\MessageBreak
       and run scripts/make-chapters.py.}}%
    {\edef\thechap{\@nameuse{chapters@#1}}}}
\makeatother
\chapnumber{<sheet ID>}
\makeatletter
\providecommand\chapters@set[2]{\@namedef{chapters@#1}{#2}}
\chapters@set{20261125T001000}{1}% zeros-to-ramp-sums
\chapters@set{20261125T001100}{2}% inert-primes-and-estimate-b
\chapters@set{20261125T005000}{3}% galois-fourier-and-averaging
\chapters@set{20261125T005500}{4}% columns-from-a-sphere
\chapters@set{20261125T000000}{5}% slope-method
\chapters@set{20261126T030100}{6}% a-twist-for-each-prime
\chapters@set{20261125T003000}{7}% the-chain-and-the-square-root
\chapters@set{20261125T010000}{8}% rank-lemma-and-zero-estimate
\chapters@set{20261125T010500}{9}% poincare-series-and-near-orthogonality
\chapters@set{20261124T223000}{10}% elliptic-curves-for-the-proof
\chapters@set{20261125T009500}{11}% fock-model-and-theta-intertwiners
\chapters@set{20261124T223100}{12}% complex-multiplication-and-inert-reduction
\chapters@set{20261125T008000}{13}% one-cm-determinant
\chapters@set{20261124T214500}{14}% supersingular-theta-divisibility
\chapters@set{20261124T214600}{15}% the-isotropic-plane
\chapters@set{20261125T013700}{16}% pairings-on-a-formal-group-of-height-two
\chapters@set{20261125T013500}{17}% hecke-nondegeneracy-on-the-tree
\chapters@set{20261126T030200}{18}% the-bound-is-attained
\chapters@set{20261125T009000}{19}% kronecker-limit-and-chowla-selberg
\chapters@set{20261125T002000}{20}% definite-theta-lifts
\chapters@set{20261125T015500}{21}% siegel-weil-in-the-smallest-case
\chapters@set{20261125T002100}{22}% definite-theta-lifts-and-norms
\chapters@set{20261125T016500}{23}% hecke-operators-and-theta-lifts
\chapters@set{20261125T016700}{24}% divisibility-along-the-chain
\chapters@set{20261125T004000}{25}% denominators-and-sha
\chapters@set{20261125T004100}{26}% tate-twists-and-sha
\chapters@set{20261126T020430}{27}% the-minus-part-by-gauss-sums
\chapters@set{20261125T014100}{28}% congruence-modules
\chapters@set{20261125T014500}{29}% eisenstein-congruences-and-selmer-classes
\chapters@set{20261125T014700}{30}% pseudocharacters-and-the-length-bound
\chapters@set{20261125T011000}{31}% galois-representations-black-box
\chapters@set{20261126T030300}{32}% the-data
\chapters@set{20261126T020210}{33}% bessel-poincare-series-and-their-gram-operators
\chapters@set{20261126T030400}{34}% the-seesaw-identity
\chapters@set{20261126T030500}{35}% divisibility-of-the-composite-lift
\chapters@set{20261126T030600}{36}% theorem-a-short-forms
\chapters@set{20261126T030700}{37}% the-determinant-of-a-chain
\chapters@set{20261126T030800}{38}% the-cost-of-one-lift
\chapters@set{20261126T030900}{39}% denominators-of-projections-to-lifts
\chapters@set{20261126T031000}{40}% twists-with-few-congruences
\chapters@set{20261126T031100}{41}% theorem-c-and-the-main-inequality
\providecommand\chapnumber[1]{%
  \@ifundefined{chapters@#1}%
    {\PackageError{chapters}{Sheet #1 is not listed in the index}%
      {List it in 20261125T012000--grh-workbook__guide.org\MessageBreak
       and run scripts/make-chapters.py.}}%
    {\edef\thechap{\@nameuse{chapters@#1}}}}
\makeatother
\chapnumber{20261125T005500}
\renewcommand{\thesection}{\thechap.\arabic{section}}

\theoremstyle{definition}
\newtheorem{problem}{Problem}[section]
\theoremstyle{plain}
\newtheorem{theorem}{Theorem}[section]

\newcommand{\Z}{\mathbb{Z}}
\newcommand{\Q}{\mathbb{Q}}
\newcommand{\R}{\mathbb{R}}
\newcommand{\C}{\mathbb{C}}
\newcommand{\PP}{\mathbb{P}}
\newcommand{\EE}{\mathbb{E}}
\newcommand{\cell}{\operatorname{cell}}

\setlength{\emergencystretch}{2em}

\title{Places from a sphere}
\author{}
\date{November 25, 2026}

\begin{document}
\maketitle
\refstepcounter{chap}\label{chap}

\paragraph{Prerequisites.}
Linearity of expectation, the union bound and Chebyshev's inequality for sums of independent Bernoulli variables.  From Chapter~\ref{gf-chap}, \emph{Fourier analysis on the Galois group}: the group \(H\), the cells \(\cell_t(r)\) and Problem~\ref{gf-pr:outside} (Section~\ref{gf-sec:places}), and the thick cells \(\Theta_u\), the density \(\delta = \lvert T_2\rvert/n\) and Theorem~\ref{gf-thm:two-scales} (Section~\ref{gf-sec:thick}).

\medskip
In Chapter~\ref{gf-chap} a set \(T_2\) of columns was fixed, and the cells of row~\(1\) in the columns \(uT_2\) were given the weight \(b\), so that they test primes near \(x\).  Nothing was asked of \(T_2\) there except that it be large: Theorem~\ref{gf-thm:two-scales} divides by \(\delta\).  How is \(T_2\) chosen?  The argument asks for a second property, that the cells of the odd powers of one prime rarely fall in \(uT_2\), and the two properties pull in opposite directions.  Behrend's construction \cite{Beh} satisfies both: lattice points on a sphere, written as digits in a large base.  We follow \cite[\S III.3]{Rub} and \cite[\S5]{Adr}.

Throughout, \(x\) is large, \(0 < \varepsilon < \frac1{52}\), and \(n\) is a power of \(2\) with \(x^{2+8\varepsilon} \le n \le 2x^{2+8\varepsilon}\); the notation \(q, d, N = 4n, H, \Phi, v(r)\) is that of Chapter~\ref{gf-chap}.  The integer \(L \ge 2\) is at most \(C_L \log x\) for a constant \(C_L\); in the argument \(L\) is of order \(\log x\) \cite[\S III.3]{Rub}.

\section{Two requirements}\label{sec:two}

\begin{problem}\label{pr:cyclic}
The group \(H = (\Z/N\Z)^\times/\{\pm1\}\) is cyclic of order \(n\), generated by \([5]\).  For \(N = 16\), \([5]^0, [5]^1, [5]^2, [5]^3\) are \([1], [5], [7], [3]\).
\end{problem}

We write \(H\) additively through \(\iota \colon H \to \Z/n\Z\), \(\iota([5]^i) = i\), and regard a set of integers \(T_2 \subset [0, n)\) as the set of columns \(\{[5]^s : s \in T_2\}\).  Then the thick cells for the translate \(u\) are the cells in row \(1\) whose column \(c\) has \(\iota(c) \in \iota(u) + T_2\).

\begin{problem}\label{pr:frob-prog}
Let \(p\) be a prime inert in \(K = \Q(\sqrt{-d})\), \(p \nmid 2qd\), and \(\iota_p := \iota([p])\).  For odd \(h \ge 1\) and \(t \in H\), the cell \(\cell_t(p^h)\) is in row \(v(p)\) and its column has index \(\iota(t) + h\iota_p\).  So the cells of \(p, p^3, p^5, \dots\) run along the progression \(\iota(t) + \iota_p + 2\iota_p\,\Z\) in a single row.
\end{problem}

We call these the \emph{Frobenius progressions}.  The first requirement on \(T_2\) comes from averaging.

\begin{problem}\label{pr:avg-density}
Let \(r\) be a prime power with a cell, and \(t \in H\).  The number of \(u \in H\) for which \(\cell_t(r)\) is thick for \(w_u\) is \(\lvert T_2\rvert\) if \(v(r) = 1\), and \(0\) if \(v(r) = 0\).
\end{problem}

So the average over translates sees \(T_2\) only through \(\delta\), and every prime power of row~\(1\) gets the weight \(\delta\); this is Problem~\ref{gf-pr:thick-avg}.  The gain of the argument is proportional to \(\delta\), while its errors are not, and in Theorem~\ref{gf-thm:two-scales} the mean excess is divided by \(\delta b\).  The argument needs
\begin{equation}\label{eq:density}
  \delta \ge x^{-\varepsilon} ,
\end{equation}
which is used in the chapter \emph{Assembly at one scale}, where the other losses are also powers \(x^{O(\varepsilon)}\).

The second requirement concerns a single translate.  For all \(u, t \in H\) and every prime as in Problem~\ref{pr:frob-prog}, at most two of the cells of \(L\) consecutive odd powers \(p^h\) may be thick, provided \(2\iota_p\) has order greater than \(4L\).  If the order is at most \(4L\), at most one cell of any odd power of \(p\) may be thick.  In \cite{Rub} the sets of columns enter the main inequality through \eqref{eq:density} and through a bound for the losses at inert primes, which is quoted there from the formal proof \cite[\S III.6]{Rub}.  These are losses of the kind in Theorem~\ref{ip-thm:main} of Chapter~\ref{ip-chap}, the weight \(w\) subtracted at each odd level, not the colength of Chapter~\ref{hn-chap}: at the scale \(p^h\) the loss is at most \(w\), where \(w\) is the weight of the cell to which Frobenius links the chosen one \cite[\S I.4.3]{Rub}, and a thick cell carries the large weight \(b\).  This second property is the form in which \(T_2\) is recorded for that bound, which we do not treat here.  The requirement is uniform: it must hold for every \(u\), \(t\) and \(p\), with no exceptions.

Why are these two requirements in tension?  The obvious sets of integers without three-term progressions are thin.

\begin{problem}\label{pr:base3}
Let \(A_k\) be the set of integers in \([0, 3^k)\) whose base-\(3\) digits are \(0\) or \(1\).  (i) \(A_k\) contains no three-term arithmetic progression, and \(\lvert A_k\rvert = (3^k)^{\log 2/\log 3}\).  (ii) For \(k \ge 2\), \(A_k\) contains \(0, 1, 3\), and \(3(1 - 0) = 1 \cdot (3 - 0)\).
\end{problem}

As a set of columns with \(n \approx 3^k\), \(A_k\) has density \(n^{-0.369\ldots}\), far below \(x^{-\varepsilon}\).  A random set has the right density but meets progressions too often.

\begin{problem}\label{pr:random-ap}
Let \(R \subset \Z/n\Z\) contain each element independently with probability \(\delta\).  The expected number of pairs \((s, h)\) with \(s, s + h, s + 2h \in R\) and \(h\) of order greater than \(4L\) is at least \(n(n - 4L)\delta^3\).  For \(\delta = x^{-\varepsilon}\) this exceeds \(\delta n \cdot x^{2}\) for large \(x\).
\end{problem}

Each such pair gives three consecutive terms \(s, s+h, s+2h\) of a progression in \(R\), against the second requirement for every \(L \ge 3\).  In expectation there are far more bad pairs than points of \(R\), so counting one deleted point for each bad pair leaves no lower bound for what survives.  The count alone does not show that no cleverer deletion succeeds, since one point can lie on many bad pairs; but it locates the density at which the naive deletion stops working.

\begin{problem}\label{pr:random-deletion}
Let \(R\) be as in Problem~\ref{pr:random-ap}, and let \(R'\) be obtained by deleting one point from each bad pair.  Then \(R'\) contains no bad pair, and its expected size is at least \(\delta n - \delta^3 n^2\).  For \(\delta \le (2n)^{-1/2}\) this is at least \(\delta n/2\), and \((2n)^{-1/2} < x^{-1-4\varepsilon}\).
\end{problem}

So a random set with deletions reaches a density about \(n^{-1/2}\), and to raise it to \(x^{-\varepsilon}\) the set must be built with structure.

\section{Behrend's construction}\label{sec:behrend}

Three points of a sphere are never collinear.  Behrend's idea is to transport this to the integers by writing vectors as digits.  Let \(S, D \ge 2\) be integers and \(\Sigma_k := \{v \in \{0, \dots, D-1\}^S : \lvert v\rvert^2 = k\}\).

\begin{problem}\label{pr:line-sphere}
Let \(v_1, v_2, v_3\) lie on one sphere \(\lvert v\rvert^2 = k\) in \(\R^S\), and let \(i, j\) be real with \(i, j \ne 0\), \(i \ne j\) and \(j(v_2 - v_1) = i(v_3 - v_1)\).  Then \(v_1 = v_2 = v_3\).
\end{problem}

\begin{problem}\label{pr:pigeonhole}
Some \(k\) has \(\lvert\Sigma_k\rvert \ge D^{S-2}/S\).
\end{problem}

Fix such a \(k\).  Let \(B := 4LD\), \(s(v) := \sum_{m=0}^{S-1} v_m B^m\), and \(T_2 := s(\Sigma_k)\).

\begin{problem}\label{pr:digits}
If \(c_0, \dots, c_{S-1}\) are integers with \(\lvert c_m\rvert < B\) and \(\sum_m c_m B^m = 0\), then every \(c_m = 0\).  In particular \(s\) is injective.
\end{problem}

\begin{problem}\label{pr:no-relation}
There are no distinct \(s_1, s_2, s_3 \in T_2\) and integers \(i \ne j\) with \(0 < \lvert i\rvert, \lvert j\rvert \le L\) and \(j(s_2 - s_1) = i(s_3 - s_1)\).
\end{problem}

For \(L = 2\) this excludes three-term progressions (\(i = 1\), \(j = 2\)).  Problem~\ref{pr:base3}(ii) shows that it says more.

\begin{problem}\label{pr:example}
For \(S = 3\), \(D = 3\), \(L = 2\), so \(B = 24\), the sphere with the most points is \(k = 5\), with the six permutations of \((0, 1, 2)\), and \(T_2 = \{26, 49, 578, 624, 1153, 1176\}\).
\end{problem}

It remains to choose \(S\) and \(D\).  Behrend lets \(S\) grow with \(n\).  For the density \(x^{-\varepsilon}\) a fixed dimension suffices: let \(S > 12/\varepsilon\) be a fixed integer and \(D := \lfloor (n/4L)^{1/S}/(4L) \rfloor\).

\begin{problem}\label{pr:size}
(i) \(B^S \le n/(4L)\), so \(4Ls < n\) for every \(s \in T_2\).  (ii) If \((n/4L)^{1/S} \ge 8L\), then \(\lvert T_2\rvert \ge n^{1 - 2/S}(8L)^{-S}/S\).
\end{problem}

Since \(n^{-2/S} \ge (2x^{2+8\varepsilon})^{-\varepsilon/6} \ge x^{-\varepsilon/2}\) and \((8L)^{-S}/S \ge (8C_L\log x)^{-S}/S \ge x^{-\varepsilon/2}\) for large \(x\), Problems~\ref{pr:pigeonhole}--\ref{pr:no-relation} and~\ref{pr:size} prove the following.

\begin{theorem}[Behrend]\label{thm:behrend}
For \(x\) large in terms of \(\varepsilon\) and \(C_L\) there is a set \(T_2\) of nonnegative integers with \(4Ls < n\) for \(s \in T_2\), with \(\lvert T_2\rvert \ge x^{-\varepsilon}n\), and with the property of Problem~\ref{pr:no-relation}.
\end{theorem}

Optimizing \(S\) gives the density \(L^{-1}\exp(-C\sqrt{\log n \log L})\) \cite[Prop.~5.17]{Adr}, which is Behrend's bound for \(L = 2\) \cite{Beh}.  No construction can give a density bounded below: by Roth's theorem \cite{Ro} a set of integers in \([0, n)\) without three-term progressions has \(o(n)\) elements.

Now pass to \(\Z/n\Z\).  The condition \(4Ls < n\) turns a congruence into an equation.

\begin{problem}\label{pr:long-order}
Let \(u_0, t, h \in \Z/n\Z\) with \(h\) of order greater than \(4L\).  Among any \(L\) consecutive terms \(t + ih\), at most two lie in \(u_0 + T_2\).
\end{problem}

\begin{problem}\label{pr:short-order}
If \(h\) has order at most \(4L\), at most one element of \(t + \Z h\) lies in \(u_0 + T_2\).
\end{problem}

Problem~\ref{pr:long-order} uses only that \(h\) has order at least \(L\); the threshold \(4L\) is the one at which Problem~\ref{pr:short-order} takes over.  Problems~\ref{pr:long-order} and~\ref{pr:short-order} prove the following.

\begin{theorem}\label{thm:progressions}
Let \(T_2\) be as in Theorem~\ref{thm:behrend}, and \(u_0, t, h \in \Z/n\Z\).  If \(h\) has order greater than \(4L\), at most two of any \(L\) consecutive terms \(t + ih\) lie in \(u_0 + T_2\).  If \(h\) has order at most \(4L\), at most one element of \(t + \Z h\) does.
\end{theorem}

So \(T_2\) satisfies~\eqref{eq:density} and the second requirement of Section~\ref{sec:two}; for the cells this reads as follows.

\begin{problem}\label{pr:frob-thick}
Let \(p\) be as in Problem~\ref{pr:frob-prog}, and \(u, t \in H\).  If \(2\iota_p\) has order greater than \(4L\) in \(\Z/n\Z\), then at most two of \(\cell_t(p^{h_0}), \cell_t(p^{h_0 + 2}), \dots, \cell_t(p^{h_0 + 2L - 2})\) are thick for \(w_u\), for every odd \(h_0 \ge 1\).
\end{problem}

Theorems~\ref{thm:behrend} and~\ref{thm:progressions} are \cite[Thm.~III.3.1]{Rub}.  For \(L(s, \psi)\) with \(\psi\) of order \(o\), Section~\ref{gf-sec:subgroup} requires \(T_2 \subset A_0\).  One applies the construction with \(n/o\) in place of \(n\) and places \(s\) at \(s\,a_0\), with \(a_0\) a generator of the cyclic group \(A_0\) \cite[\S5]{Adr}.

\section{The tops: a random choice}\label{sec:tops}

The set \(T\) of shifts \(t\) in Section~\ref{gf-sec:subgroup}, the \emph{tops}, has the opposite needs.  Its density \(\delta_1 := x^{-\varepsilon/8}\) is larger, and along a Frobenius orbit it may contain a bounded number \(k_T\) of points in a window, not two.  Exceptions are allowed if they are rare on average over primes.  A random set does this, and this is where the argument uses randomness \cite[\S III.3]{Rub}, \cite[Prop.~5.20]{Adr}.  Let \(T \subset \Z/n\Z\) contain each element independently with probability \(\delta_1\).

\begin{problem}\label{pr:size-T}
\(\PP\bigl(\bigl\lvert\lvert T\rvert - \delta_1 n\bigr\rvert > \frac12\delta_1 n\bigr) \le 4/(\delta_1 n)\).
\end{problem}

\begin{problem}\label{pr:window}
Let \(h\) have order at least \(L_2\), and let \(W\) be a set of \(L_2\) consecutive terms of \(t + \Z h\).  Then \(\PP(\lvert W \cap T\rvert \ge k_T + 1) \le (L_2\delta_1)^{k_T+1}\).
\end{problem}

\begin{problem}\label{pr:union}
Let \(k_T \ge 8(4 + 8\varepsilon)/\varepsilon\).  If \(h\) has order at least \(L_2\), the expected number of windows of \(L_2\) consecutive terms on the orbits of \(h\) that contain more than \(k_T\) points of \(T\) is at most \(2L_2^{k_T+1}x^{-2-\varepsilon/8}\).
\end{problem}

With \(L_2\) bounded and \(k_T\) as in Problem~\ref{pr:union}, a random \(T\) is thus sparse along the orbits of one \(h\) with probability close to \(1\).  The argument asks for two more conditions of even spread along the orbits.  It conditions on \(t \in T\) and weights the primes by \(\log p/p\), and it takes different probabilities on the cosets of \(A_0\) to obtain the shares of Problem~\ref{gf-pr:shares}.  Then it chooses one \(T\) by comparing expectations \cite[Thm.~III.3.2]{Rub}, \cite[Prop.~5.20]{Adr}.  Compare Problem~\ref{pr:random-ap}: the same method fails for \(T_2\), where at most two points are allowed in a window and there may be no exceptions.  For \(T_2\) the sphere replaces chance.  This ends the reduction: the Riemann hypothesis now asks for an upper bound for weighted prime sums at one scale, and the next part, from Chapter~\ref{sm-chap} on, gives the method that bounds them by comparing divisibility with size.

\subsection*{Exercises}

\begin{enumerate}
\item Let \(\log L \le c\log n\) with \(c\) small.  With \(S \approx \sqrt{\log n/\log L}\) and \(D := \lfloor (n/4L)^{1/S}/(4L) \rfloor\), the construction gives \(\lvert T_2\rvert \ge n\exp(-C\sqrt{\log n \log L})\), still with \(4Ls < n\), for an absolute constant \(C\).
\item Without the condition \(4Ls < n\), Problem~\ref{pr:long-order} fails: for \(n = 64\) and \(L = 3\), the set \(\{0, 2, 33\}\) has the property of Problem~\ref{pr:no-relation}, yet \(0, 33, 66 \equiv 2\) are consecutive terms of \(33i\), of order \(64\).
\item A set \(T_2 \subset [0, n)\) with the property of Problem~\ref{pr:no-relation} meets every arithmetic progression of \(L\) terms in \(\Z\) in at most two points.
\end{enumerate}

\section*{Hints}

\subsection*{\S\ref{sec:two}}
\begin{description}
\item[\ref{pr:cyclic}] By induction, \(5^{2^{e}} \equiv 1 + 2^{e+2} \pmod{2^{e+3}}\); and \(-1\) is not a power of \(5\) modulo \(4\).
\item[\ref{pr:frob-prog}] Problem~\ref{gf-pr:outside}; \(\chi_q(p^h) = \chi_q(p)\) for odd \(h\).
\item[\ref{pr:avg-density}] Count the solutions \(u\) of \(t[r] \in uT_2\).
\item[\ref{pr:base3}] (i) Add \(a + c\) digit by digit.
\item[\ref{pr:random-ap}] The three points are distinct; count the \(h\) of order at most \(4L\) in a cyclic group.
\item[\ref{pr:random-deletion}] Linearity of expectation: \(\lvert R'\rvert\) is at least \(\lvert R\rvert\) minus the number of bad pairs.
\end{description}

\subsection*{\S\ref{sec:behrend}}
\begin{description}
\item[\ref{pr:line-sphere}] Expand \(\lvert v_1 + b\rvert^2 = \lvert v_1\rvert^2\) for \(b = v_2 - v_1\) and for \(c = v_3 - v_1\), and eliminate \(v_1 \cdot c\).
\item[\ref{pr:pigeonhole}] Count the possible values of \(\lvert v\rvert^2\).
\item[\ref{pr:digits}] Look at the lowest nonzero \(c_m\) modulo \(B\).
\item[\ref{pr:no-relation}] Apply Problem~\ref{pr:digits} to the digits of \(j(s_2 - s_1) - i(s_3 - s_1)\).
\item[\ref{pr:size}] (ii) \(\lfloor z\rfloor \ge z/2\) for \(z \ge 1\), and Problem~\ref{pr:pigeonhole}.
\item[\ref{pr:long-order}] Eliminate \(h\) from three such terms, and bound the resulting integer.
\item[\ref{pr:short-order}] If \(m\) is the order of \(h\), then \(m(s_1 - s_2) \equiv 0\).
\item[\ref{pr:frob-thick}] Problems~\ref{pr:frob-prog} and~\ref{pr:long-order}.
\end{description}

\subsection*{\S\ref{sec:tops}}
\begin{description}
\item[\ref{pr:size-T}] Chebyshev's inequality.
\item[\ref{pr:window}] Union bound over subsets of \(W\) of size \(k_T + 1\).
\item[\ref{pr:union}] There are \(n\) windows.
\end{description}

\begin{thebibliography}{Adr}

\bibitem[Adr]{Adr}
J.~Adriazola, \emph{The Alp\"oge proof of the Riemann hypothesis}, manuscript, revised 24 November 2026, \url{https://github.com/jimjamzola/RH-Alpoge-Reconstruction}.

\bibitem[Beh]{Beh}
F.~A.~Behrend, \emph{On sets of integers which contain no three terms in arithmetical progression}, Proc. Nat. Acad. Sci. U.S.A. \textbf{32} (1946), 331--332.

\bibitem[Ro]{Ro}
K.~F.~Roth, \emph{On certain sets of integers}, J. London Math. Soc. \textbf{28} (1953), 104--109.

\bibitem[Rub]{Rub}
\emph{The Alp\"oge--Claude proof of the Riemann hypothesis explained for the Riemann zeta function}, written by AI models at the behest of M.~O.~Rubinstein, working draft, 21 November 2026, \url{https://www.math.uwaterloo.ca/~mrubinst/alpoge_claude_grh/zeta-paper.pdf}.

\end{thebibliography}

\end{document}
